\newcommand{\legendone}{\textbf{Figure 1 | Measuring human-aligned gloss judgments and their internal correlates.}
\textbf{a,} The shared stimulus domain comprises 3,888 rendered images (36 shapes, 36 illuminations and three viewpoints). Published human matches provide an external behavioral reference for two frozen vision--language models.
\textbf{b,} Three measurement interfaces: single-image rating, two-image choice and adaptive reference matching. Prompt excerpts summarize the requested judgment; complete acquisition strings and example records are provided in Supplementary Fig.~S11. Matching uses a reconstructed 52-level reference library and at most seven comparisons per search. The controller updates the reference, not the model weights. Two starts yield two endpoints per image.
\textbf{c,} Behavioral analyses distinguish human agreement beyond physical ordering, shared condition-response geometry and prediction across tasks. Continuous task mappings are evaluated in held-out shape or illumination groups. Human scores enter the comparison analyses but are not supplied during behavioral acquisition.
\textbf{d,} Internal measurements separate a supervised readout of frozen states from reference-defined perturbation. Readouts use human labels within training folds; the perturbation direction is derived from the reference sequence without those labels. Signed displacements to test-image bridge tokens test influence on gloss choices while pixels and model parameters remain fixed. Figures~2--6 report the corresponding evidence. Detailed acquisition, geometry, readout and intervention workflows are retained in Supplementary Figs.~S12--S15.}

\newcommand{\legendtwo}{\textbf{Figure 2 | Model judgments track human appearance beyond physical gloss.}
\textbf{a--c,} Image-density plots compare physical gloss with human matches (a), and human matches with InternVL (b) and Qwen3 (c) reference matches; $n=3,849$ images. Model values average two searches. Dotted lines denote numerical equality, without cross-system calibration.
\textbf{d,} Spearman correlations with human judgments for ratings and matches; physical gloss provides a benchmark.
\textbf{e,} Distribution of within-level model--human Spearman correlations across all 50 exact physical-gloss levels; each level contributes one coefficient.
\textbf{f,} Human-order agreement for 1,000 random and 1,000 conflict pairs, combining both display orders and assigning half credit to ties. Human ordering is derived from mean matches; the dotted reference denotes 50\%.
\textbf{g,} Stimuli organized by human-match tertiles and equal-width physical intervals. Counts and median model percentile ranks summarize all images in each cell; displayed examples lie near cell centers.
\textbf{h,} Pearson correlations after subtracting each system's mean at each physical level.
\textbf{i,} A selected conflict example on which both models follow human ordering in both orders; aggregate evidence is in f.
Points in d, f and h are observed estimates; error bars are 95\% percentile intervals from 2,000 bootstrap samples. Resampling crosses shape and illumination in d and h and operates on image endpoints in f. Physical-level means are re-estimated in h. All matching endpoints are retained. Teal denotes InternVL and orange Qwen3. Example-selection rules and exact statistics are documented in Supplementary Note~S20; uncertainty concerns stimuli, not human-participant sampling.}

\newcommand{\legendthree}{\textbf{Figure 3 | Human and model judgments share scene-dependent organization.}
Analyses use all 3,888 images.
\textbf{a,b,} Adjusted effects for 36 shapes (a) and 36 illuminations (b), ordered by the human effect. Separate additive regressions control physical gloss, categorical chroma, hue and viewpoint; effects are in each system's response standard deviations.
\textbf{c,} Spearman correspondence of human and model condition effects. Bars show 95\% condition-bootstrap intervals (2,000 samples). Two-sided condition-label tests use 99,999 permutations; four-test Benjamini--Hochberg correction yields $q=0.00177$ for InternVL shape and $q=1.33\times10^{-5}$ otherwise.
\textbf{d,e,} Shape-profile (d) and illumination-profile (e) dissimilarity matrices after cross-view physical/color adjustment. Mean squared profile differences retain additive effects. Each matrix is divided by its mean off-diagonal distance for display, with colors saturating at three; condition order is identical across systems.
\textbf{f,} Fractional ranks of all 630 condition-pair distances in each human--model comparison. Each point is a condition pair, not an independent stimulus.
\textbf{g,} RSA across all system pairings. Points are Spearman estimates and bars 95\% node-bootstrap intervals (2,000 samples; 36 nodes). Tests permute matrix rows and columns together (4,999 permutations); all six adjusted-profile comparisons have $q\leq0.0009$ within the 18-test family (Supplementary Table~S6).
\textbf{h,} Rendered examples from conditions with extreme human-adjusted effects, selected using human estimates alone. Reported effects use the full dataset. The gallery illustrates those conditions rather than a controlled image manipulation. Supplementary Note~S20 specifies example selection and matrix-display conventions.}

\newcommand{\legendfour}{\textbf{Figure 4 | A common ordering transfers across response interfaces.}
\textbf{a,} Prediction between each model's behavioral outputs, without human calibration targets.
\textbf{b,c,} Shape-held-out predictions of reference matches from ratings (b) and ratings from matches (c), using nondecreasing isotonic regression. Each plot contains 3,849 out-of-fold predictions; color denotes image count and dotted identity lines indicate exact prediction.
\textbf{d,e,} Pooled $R^2$ for the two mapping directions under separate 12-fold shape and illumination splits. Points show observed scores; bars show 95\% intervals from 2,000 crossed-condition bootstrap samples of fixed predictions.
\textbf{f,} Agreement between rating- or match-derived ordering and independently collected pairwise choices across 2,000 pairs, retaining both presentation orders. Ties receive half credit.
\textbf{g,} Empirical choice frequency against signed rating or matching differences, normalized by their scale ranges. Lines join the means of ten fixed bins, with symbol area increasing with bin count. These descriptive summaries have no confidence bands and are not fitted psychometric curves.
\textbf{h,} Directional agreement within the 1,000 random and 1,000 human--physical conflict pairs. In f and h, points show agreement and bars show 95\% intervals from 2,000 image-node bootstrap samples over 2,419 unique endpoints.
\textbf{i,} Two concordant examples connecting the recorded tasks. Examples were selected near the lower and upper quartiles of InternVL matching differences within the random stratum; aggregate analyses retain all pairs. Teal denotes InternVL and orange Qwen3. Fold definitions, uncertainty units and example-selection rules are given in Methods and Supplementary Notes S6 and S20.}

\newcommand{\legendfive}{\textbf{Figure 5 | Frozen internal states contain information predictive of human gloss judgments.}
\textbf{a,} Readouts across visual blocks, the bridge and prompt-conditioned language blocks. Model weights remain frozen; separate ridge readouts use training labels and group-disjoint regularization selection.
\textbf{b,c,} Human-target prediction across sampled states under shape (b) and illumination (c) holdout.
\textbf{d,e,} Prediction combining a training-fitted physical/color baseline with an internal-state readout of human residuals. Scores evaluate the combined prediction against original human responses; dashed lines denote baseline-only $R^2$.
\textbf{f,} Held-out $R^2$ for human matches, physical gloss and the model's reference matches at all nine common states, plus Qwen3's three DeepStack outputs. S and I indicate shape and illumination holdout. Targets are fitted separately.
\textbf{g,} All 3,888 shape-held-out predictions from L36, reported by architectural position. Hexagon color denotes image count; dotted identity lines indicate exact prediction. Predictions are neither clipped nor rescaled.
\textbf{h,} Paired L36 gain over the physical/color baseline, $\Delta R^2=R^2_{\mathrm{combined}}-R^2_{\mathrm{baseline}}$.
Curves and points show observed scores. Shading and bars indicate pointwise 95\% intervals from 2,000 crossed shape--illumination bootstrap samples, with predictors fixed and paired weights shared in h. Each analysis uses all 3,888 images in separate 12-fold grouped evaluations. Visual/bridge states are token means; language states use the final prompt token before answer generation. Pooling differences preclude interpreting the curves as a causal hierarchy. Violet denotes InternVL and teal Qwen3. Additional implementation details and exact gains appear in Supplementary Notes S13 and S20.}

\newcommand{\legendsix}{\textbf{Figure 6 | A reference-defined visual direction bidirectionally shifts gloss choices.}
\textbf{a,b,} Direction construction from 52 reconstructed references spanning nominal gloss $c=0$--0.1487 (a), and projections of those same references (b), rescaled within model. The normalized multivariate slope defines direction $\mathbf d$ and its norm the dose scale $s$; human test labels are unused.
\textbf{c,} A selected rendered image and native bridge-token projections onto $\mathbf d$, separately color-scaled. These illustrate responses, not causal localization; selection favored visible specular structure and local-brightness association.
\textbf{d,} Uniform token perturbation $\mathbf z'_{i,t}=\mathbf z_{i,t}+\alpha s\mathbf d$. Pixels, reference tokens and model weights are fixed; Qwen3 DeepStack paths remain unchanged.
\textbf{e,f,} Baseline-subtracted test-glossier scores for InternVL (e) and Qwen3 (f). Scores normalize logits over A/B/SAME and assign half weight to SAME. Lines join empirical means; the reference-direction band is a 95\% shape-bootstrap interval. Controls comprise equal-norm luminance-associated and three orthogonal random directions.
\textbf{g,h,} Reference-minus-control differences in negative-to-positive score change. Dots denote 36 shapes; diamonds and bars show means and paired 95\% intervals. Two-sided sign-randomization tests (99,999 draws; eight-test Benjamini--Hochberg correction) yield $p=q=10^{-5}$ for all controls.
\textbf{i,} Complete shape-level dose profiles in original identifier order.
The assay uses 72 images: 36 shapes, two viewpoints, one illumination, reference 27 and both display orders. Inference averages viewpoints and orders within shape, with 2,000 shape-bootstrap samples. Effects are percentage-point changes in a candidate-normalized score, not accuracy gains. Selection rules and full assay definitions are given in Supplementary Notes S14, S18 and S20.}
